% !TEX root = ../main.tex
\chapter{Nulls and statistics}\label{app:nulls}
This appendix says how the app tells a pattern from chance: what each null is, which code
draws it, how its randomness is fixed, where its result is shown, and how to read the numbers.
\Cref{sec:core-ideas-nulls} gives the short version; the surrogate block's own entry is
\blockentry{preprocessing.surrogate}.

\section{What a null is for}\label{sec:nulls-what}
\begin{plainterms}[Would a random deck do as well?]
Any long, noisy recording contains look-alikes by accident, the way a long enough run of coin
tosses contains ten heads in a row. So ``the detector found 40 events'' says little by itself.
The useful question is: how many would it have found in a recording that has the same general
character but no real events? A \term{null}{null} is the app's answer to that question, made
by running the very same analysis on fake copies of the data, the
\term{surrogate}{surrogates}, which keep what is boring about the signal and destroy what might
be interesting. A result is evidence of something only if it beats its null.
\end{plainterms}

What a null keeps and what it destroys decides what it tests. A surrogate that keeps the
signal's frequency content but scrambles every local shape asks ``would this fire as often on
noise of the same colour?''. One that keeps short pieces intact and shuffles their order asks
about timing, not shape. One that keeps the exact values asks about everything except the
value distribution. Choosing the null is therefore part of the question, and the app prints
which null judged every number it shows.

Three rules hold everywhere in the app:
\begin{itemize}
\item \emph{The null named is the null drawn.} Every surface prints the number of draws
  actually made and the method, not a setting it might have used.
\item \emph{A surrogate detection is never a detection.} Runs made on surrogates are stored
  (linked to the real run by \code{runs.surrogate_of_run_id}) but every reader goes through
  one predicate, \code{not_surrogate} in \code{Working/database/queries.py}: they are never
  counted in Explore, sent to Review, promoted to the Library or classified across channels.
\item \emph{No null is not a pass.} Where a page draws no null, it says so in words, and a
  number there is a description, not evidence.
\end{itemize}

\section{The surrogates the app can make}\label{sec:nulls-surrogates}
\subsection{Surrogates of a whole span}
\blockref{preprocessing.surrogate} (\code{Adapters/preprocessing_surrogate.py}) turns a span
$x_0,\dots,x_{n-1}$ into a surrogate of the same length. It offers three methods; Settings ›
Nulls reads the list off the block (\code{offered_methods} in
\code{Working/discovery/seeded_search.py}), so it cannot offer one the block does not
implement.

\begin{description}
\item[Phase randomisation] (\code{_phase_randomise}; the default). Take the real Fourier
  transform $X_k$ of the span, multiply each bin by $e^{i\phi_k}$ with $\phi_k$ uniform on
  $[0, 2\pi)$, and invert. The DC bin, and the Nyquist bin for an even $n$, keep phase~0 so the
  result is real. The magnitudes $|X_k|$, and so the power spectrum, are kept exactly; every
  local shape is scrambled. The result is close to Gaussian whatever the original
  distribution. This is the null for a detection chain: \emph{would this chain fire as often
  on a signal with the same spectrum and no real events?}
\item[AAFT] (\code{aaft_surrogate}, the amplitude-adjusted Fourier transform of
  \citet{theiler1992testing}). Rank-map $x$ onto sorted Gaussian deviates, phase-randomise
  that, and rank-map the result back onto $x$'s own sorted values. The output is a
  permutation of the input's exact values, with approximately (not exactly) its spectrum;
  the iterative refinement, IAAFT, is not offered. Use it for a measure that reads the value
  distribution: \blockref{catalogue.multiscale_entropy} says on its face to judge a curve
  against AAFT copies, because plain phase randomisation makes the surrogate Gaussian and a
  skewed but structureless signal then ``beats'' it on its distribution alone (the module's
  docstring records 20 in 20). \blockref{detection.pan_matrix_profile} draws AAFT copies
  internally by default (3 draws), for the same reason.
\item[Block shuffle] (\code{_block_shuffle}). Cut the span into whole blocks of
  \emph{block\_s} seconds and permute them; a trailing part-block stays in place. Any motif
  shorter than a block survives intact, so this null tests timing and order (trains,
  intervals), not shape. A block shorter than two samples is refused (it would be a sample
  shuffle, which is neither null), and so is an unset block length. When the app pairs a block
  shuffle with a detection run, the block is twice the longest span the real run found
  (\code{surrogate_draw_params} in \code{Working/run_groups.py}); if the run found nothing,
  no null is drawn and the reason is printed. For a seed search it is twice the seed. Settings
  › Nulls can fix one length for all.
\end{description}

\begin{watchout}[What each surrogate still keeps]
Phase randomisation keeps the spectrum, so a strong rhythm in the recording is partly in the
null too: a periodic signal's period survives. The block's view of
\blockref{detection.pan_matrix_profile} says this in its rules. Block shuffle keeps every
shape shorter than a block, so it can never show that a shape is unusual.
\end{watchout}

\subsection{Nulls made from windows rather than whole spans}
Some questions are about a set of short windows (a Library family, a stack of cycles) rather
than a span. Two further nulls serve them.
\begin{description}
\item[Trend-kept surrogates] (\code{trend_kept_surrogates} in
  \code{Working/library/grouping/methods/kshape.py}). Per window, fit a cubic, phase-randomise
  what is left and add the cubic back. A window of strongly coloured noise is mostly a slow
  curve, and plain phase randomisation turns that curve into a wiggle, so groups of noise
  crowd less than real noise does and a test against them flatters noise. Keeping the cubic
  fixes that; the degree was chosen by measurement on synthetic noise (the comment above
  \code{TREND_DEGREE} records it). Used by the Library's HDBSCAN and k-Shape noise tests and by
  the coherence score (\code{Working/library/coherence.py}), 5 draws.
\item[Quiet-stretch windows] (\code{Working/background_null.py}). Windows cut at random from
  the recording's own stretches where no event lies and no sample is missing, at the lengths
  of the real windows. They carry the recording's colour, drift and artefacts and none of its
  shapes, so a smooth shape that fills its window, which reads ``noise-like'' against
  trend-kept surrogates, can beat them. For a Library family each member gets a window from
  its own recording; the quiet-stretch null judges only when at least half the members got
  one, otherwise the trend-kept null judges and the page says so. Since 15~October 2026 a
  family passes if it beats \emph{either} null. \blockref{interrogation.cycle_shapes} uses the
  same sampler for its pseudo-cycles, and \blockref{interrogation.event_families} for its
  families.
\end{description}

\subsection{Nulls that are not surrogates of the signal}
\begin{description}
\item[Random times, across channels] (\code{chance_null} in \code{Working/cross_channel.py}).
  A member and a sibling channel count as a pair only if their correlation beats the same
  sibling at 100 random other times (windows of the member's length, at least 60\,s away,
  never over a human-marked artifact), above the 95th percentile. The draw is seeded from the
  member id. Both numbers are settings (Analysis defaults).
\item[Label shuffle, for models] (\code{Working/training/paired.py}). The training labels are
  permuted and the forest refitted, 200 times per arm by default. The one-sided p-value is
  $(1 + \#\{\text{shuffles} \geq \text{observed}\})/(1 + n)$ (\code{null_p} in
  \code{Working/training/metrics.py}), so with 200 shuffles the smallest possible p is
  $1/201 \approx 0.005$. That is why the baseline run reports $p = 0.005$ for both arms: it is
  the floor, not a measurement of how far above chance they are.
\end{description}

\subsection{How the randomness is fixed}
Every surrogate is drawn from a seeded generator, never from numpy's shared global state.
\blockref{preprocessing.surrogate} takes its \emph{seed} as a parameter, so the seed is part
of the step's recipe and the recipe's hash reproduces the surrogate exactly. A paired null of
$N$ draws runs draw $i$ at seed $s+i$ (\code{run_paired_recipe} and \code{draw_paired_null}
in \code{Working/run_groups.py}); the bridge passes base seed $s = 0$ for both Analyse and
Discovery runs. Each draw is its own recipe, so a draw already made is reused by hash:
raising the count from 20 to 50 adds 30 draws and re-runs none. The draws a run made are
linked to it in the database, and the run reports how many it made, fewer than asked if it was
cancelled. The seed search's own null (\code{null_distances}) follows the same rule, seed
$s+i$ for draw $i$. The window nulls above take a \emph{seed} parameter on their blocks.

\begin{watchout}[A fixed seed is one answer, not the right one]
Because the base seed is always 0, every run sees the same surrogates. That makes results
reproducible, but a count that depends on the particular draws will not show it. Where a
number matters, re-run with a different seed from a script and see whether it moves; the
symbolic encoders' fixed seeds were checked this way (\code{docs/rq_roundA/RQ4-bands-and-symbols.md}).
\end{watchout}

\section{Null sweeps}\label{sec:nulls-sweeps}
\begin{plainterms}[Testing every setting against chance]
A null sweep asks the null question at several values of one setting at once: at each value,
how many does the real signal give, and how many does the null give? The setting where the
real count pulls clear of the null is a principled choice, rather than a guess.
\end{plainterms}

The Block page (\apppage{analyse.block}) was designed with a null-sweep panel. On the live
page that panel is a card, \emph{This parameter against the null}, which says that a null is
drawn per chain run, not per block, and points to the toolbar. The sweep charts with values to
click exist only in the demonstration chains, which run on example data (the components under
\code{webui/client/src/analyse/demo/}); no live block page sweeps a parameter against its null.

What the live app does instead:
\begin{itemize}
\item \textbf{A null per run, in Analyse.} Turn on the toolbar's \ui{surrogate null}. After the
  real run, the whole chain runs again on surrogates of the span, with the method and draw
  count of Settings › Nulls for detection chains (20 phase-randomised draws by default). The
  footer reads, for example, ``detected 120 · null expects 31 over 20 phase randomize draws ·
  3.9× null''. To compare settings, change one, run again, and compare the footers.
\item \textbf{A null per channel, in Discovery.} Every template run in
  \apppage{discovery.runs} carries its paired null. The scoreboard's \emph{null expects} is
  the mean count over the draws on that channel, and \emph{× null} is the real count divided
  by it (\code{channel_score} in \code{Working/discovery/scoreboard.py}). A run total sums the
  channels that carry a null and says how many do.
\item \textbf{The seed search's histogram.} \apppage{discovery.seed} draws the distance of each
  real match with the null's matches behind it, and the cut can be dragged: at any cut the page
  reads how many matches it keeps and how many the null gives per draw. This is a sweep over
  the cut, read live from draws already made (\code{cut_counts}). With a scale bank, each
  length has its own null.
\item \textbf{Nulls inside a block.} Some blocks judge their own output and print the verdict:
  \blockref{detection.pan_matrix_profile} (which window lengths repeat beyond AAFT noise),
  \blockref{interrogation.cycle_shapes}, \blockref{interrogation.event_families} and
  \blockref{catalogue.shape_families} (each group against quiet stretches or trend-kept
  surrogates). The noise floor of the symbolic encoder \blockref{detection.sax_dsax} is set
  from surrogates of the span, and the count it draws grows as the span shrinks.
\end{itemize}

\screenshot{pages/zzz_surrogates/analyse.chain--fixup-t-run-with-null-reads-detected-versus-surrogate}{An
Analyse run with the surrogate null on: the footer reads what was detected against what the
null expects, and how many draws it made.}

\subsection{The Nulls settings page}
\apppage{settings.nulls} is where the null is chosen for the project. What the core reads from
it, traced through \code{null_from_settings} and \code{cut_rule_from_settings} in
\code{Working/discovery/seeded_search.py}:
\begin{itemize}
\item the method and draw count for \emph{Detection chains} (\code{null.detection.*}, default
  phase randomisation, 20 draws) and for \emph{Seed search} (\code{null.seed-search.*}, 200
  draws). A detection draw is the whole chain again, hence the smaller default;
\item \emph{significance level α} (default 0.01) and \emph{multiple channels}, the correction
  (none, Holm or Benjamini--Hochberg; default none). These reach the seed search's recommended
  cut only;
\item \emph{block shuffle · block length} (0 means twice the longest motif under test).
\end{itemize}
The page also lists rows for interrogation distributions and intervals, training baselines,
full-model shuffles and Library groupings, a \emph{seed} column, \emph{show × null beside
counts} and \emph{reuse null draws while the recipe is unchanged}. Those values are saved, but
no code in \code{Working/} or \code{webui/server/} reads them: the training job, for example,
takes its 200 shuffles from its own recipe, and draws are reused by recipe hash whatever the
toggle says. The page warns when α is smaller than $1/N$ for the smallest draw count among the
kinds that report a p-value, since no smaller p can be reported.

\section{Reading the numbers}\label{sec:nulls-reading}
\begin{plainterms}[What ``beats chance'' does and does not mean]
Beating a null says only that the result is unlikely to be produced by the kind of noise that
null makes. It does not say the result is biologically meaningful, that the detector found the
right events, or that a different null would agree. A result that beats a weak null (one that
destroys too much) proves little; a result that fails a strong one may still be real but
unresolvable with this data.
\end{plainterms}

\subsection{× null and ``null expects''}
For detection runs the app reports a count and a ratio, never a p-value. With $N$ draws
finding $c_1,\dots,c_N$ events, \emph{null expects} is $\bar c = \frac1N\sum c_i$ and
\emph{× null} is the real count divided by $\bar c$. A ratio of 3 means three times as many
events as the noise gave on average. It says nothing about how variable the null counts were:
with 20 draws, a ratio of 1.3 can be well within the spread. When the null found nothing the
ratio is undefined, and the page says ``the null found nothing'' rather than printing
infinity.

\subsection{The seed search's p and cut}
For each channel the seed search takes its closest real match and counts how many null matches
per draw lie at or below that distance; that rate, capped at 1, is the channel's $p$
(\code{recommended_cut_corrected}). Strictly it is an expected count per draw, not a
probability, and the two agree only when it is small. The recommended cut is the largest real
distance at which the null still gives at most α matches per draw: with 200 draws and
α~=~0.01, the null may contribute at most two matches below the cut over all 200 draws
(\code{recommended_cut}). When even the closest match is no better than chance, there is no
recommended cut, and the page says so.

\subsection{Model p-values}
A label-shuffle p-value (above) asks whether a model learned anything at all from its labels.
It does not compare two models; the paired difference, its bootstrap interval and McNemar's
test do that (\apppage{models.compare}). The interval of the difference is the number to read:
a difference counts when its 95\% interval excludes zero.

\subsection{z-scores and ratios to noise}
Several blocks report how far a result sits from its null in units of the null's spread, or as
a ratio to a noise score: the pan matrix profile's gap to the null, a family's coherence
against its trend-kept or quiet null, the cross-channel percentile test. Each block prints the
rule beside its number, and that rule, not the general idea of a z-score, is what the number
means. A distance or an entropy cannot go below zero and is not normally distributed, so read
these as the block defines them, not as a standard normal deviate.

\subsection{Many comparisons at once}
Testing many channels, bands or lengths at once multiplies the chances of a lucky result:
at α~=~0.01, one test in a hundred passes by chance, so across sixteen channels and three bands
about one false pass in two runs is unsurprising. The app offers two corrections for the seed
search's per-channel cut. \emph{Holm} \citep{holm1979simple} controls the chance of
\emph{any} false pass: it holds the smallest $p$ to $\alpha/m$, the next to $\alpha/(m-1)$, and
stops at the first failure. \emph{Benjamini--Hochberg} \citep{benjamini1995controlling}
controls the expected \emph{share} of false passes among those passed: with $k$ the largest
rank whose $p_{(k)} \leq \alpha k/m$, the $k$ smallest pass. Both are applied across the
channels in scope (\code{corrected_levels}), and the sentence beside the cut names the one
used. Nothing else in the app corrects for multiple comparisons: a Discovery run over sixteen
channels, a band application over several bands, or a scale bank over three lengths reports
each × null on its own. Correct for them yourself when you report, and say which correction
you used.

\begin{watchout}[A null does not travel]
A null is drawn for one scope: one span, one channel, one band, one length. Do not carry a
null count from one scope to another, from a short preview to a long run, or from one block's
settings to another's. The app keeps a separate null per channel, band and scale length for
this reason.
\end{watchout}
